\documentclass[11pt]{article}
\usepackage[T1]{fontenc}
\usepackage{lmodern,microtype}
\usepackage{amsmath,amssymb,booktabs,enumitem}
\usepackage[letterpaper,margin=1in]{geometry}
\usepackage[colorlinks=true,urlcolor=blue,citecolor=blue]{hyperref}
\newcommand{\F}{\mathbb F}
\newcommand{\Z}{\mathbb Z}
\title{Metaplectic Jacquet--Langlands Transfer, Companion Paper (Rank Two, Wild Place)\\Revision Log: Draft 6 to Draft 7}
\author{Prepared by Claude (Anthropic) for David Joyner}
\date{November 8, 2026}
\begin{document}
\maketitle
\section*{Scope}
Draft 7 responds to ChatGPT's referee report on Draft 6 (seven
recommendations, all adopted), adds the material on representation
levels requested by David Joyner, and, following his decision to stop
the research at level two, adds a Future Directions section (new
Section 15) stating what a general result would need.  No new theorem.

\section*{Decision to stop at level two}
Recorded in Section 15.1.  The reasons: levels are unbounded, so no
finite list proves Target B; the level-two proof uses three features
that all fail at level three (class-two group with derived subgroup of
order $3$; $\Theta_\rho\circ r$ supported on a central subgroup; a
central subgroup $\hat A$ large enough to detect every $t\notin Q_3$ through
central characters); and the referee's analysis agrees that the central
problem is now an inductive mechanism, not a computation.  Target B is
therefore stated as: proved for every row of level $\le2$, open beyond.

\section*{The referee's seven corrections}
\begin{enumerate}[leftmargin=*]
\item \textbf{Section 13.5 (now ``Beyond level two'').}  The false
statement that $Q_4/Q_{3m+1}$ is abelian is removed ($[1+\varpi^2i,1+\varpi^2j]$
has depth $9$, so $Q_4/Q_{10}$ is nonabelian); the candidate abelian
normal subgroup is $Q_{\lceil(3m+1)/2\rceil}/Q_{3m+1}$, i.e.\ $Q_5/Q_{10}$
at $m=3$.  The sentence on linear $\chi$ now refers to $D_8$-invariant
sums.  The ``single free orbit'' statement is corrected: the $243$
genuine characters of $\hat H_1=\tilde V$ form nine free orbits, of which
$\Lambda$ is one.  $H_m$ and $\hat H_m$ are standardized with the cutoff
$N_m=\max(3m+1,6)$.  The two statements (B1), (B2) are formulated
separately, as recommended, with the remark that a failure of (B1) would
be a genuine obstruction while a failure of universal freeness would not.
The referee's explicit non-integrality example
($\chi(a+bi)=\zeta^a$, $h=(0,0,0,1,1)$, inner product $1/3$) was
reproduced from the tables and is quoted.
\item \textbf{Howe-duality remark (now Remark 16.1).}  The pointwise
kernel $K_\theta$ diverges ($K_\theta(1,1)=27\sum\dim\rho^2$); the remark
now works with the finite truncations $K^{(m)}_\theta$ over the
irreducible representations of $H/K^{m+1}$ with central character
$\theta$, which are honest class functions on $\tilde H\times H/K^{m+1}$
with the explicit formula, and Target B is restated as ``$K^{(m)}_\theta$
is a character for every $m$''.  The Schr\"odinger-model language in
Section 2.3 is now explicitly an analogy: $\tilde U_3/N_6$ is abelian
and $\tilde U_1/N_6$ is not a Heisenberg group with nondegenerate
symplectic quotient; what is used is Clifford theory for an abelian
normal subgroup.
\item \textbf{Self-contained level-two theorem.}  New Lemma 13.7 (the
group $H=Q_3/Q_7$): $H'=Q_6/Q_7$, $Z(H)=Q_4/Q_7$, $243+54$ characters,
and $3\sum_{\chi'(1)=3}\chi'=\mathrm{Reg}_H-\mathrm{Inf}\,\mathrm{Reg}_{H/H'}$
together with the $\omega$-isotypic form, proved by hand; Step 2 of the
theorem now cites it.
\item \textbf{$\beta\neq0$ made explicit} (Lemma 10.2(b)).  For
$E=F(j)$: $(1+aj)^3=1+a^3\varpi j+3a^2\varpi+3aj$ with $v_E(3)=4$, so
cubes of $U_E^{(1)}$ have leading term at $E$-depth $3$ and their
$E$-depth-$5$ component is determined; hence $E^\times/F^\times E^{\times3}$
(order $9$) has basis the classes of $1+j$ and $1+\varpi^2j$, the
commutator pairing is nondegenerate on it, and
$[\widetilde{1+j},\widetilde{1+\varpi^2j}]=\kappa_{1,5}(1,1)=\zeta^{2\beta}\neq1$.
The consistency with the top-level Hilbert pairing of $E$ at level $6=pe_E/(p-1)$
is noted.
\item \textbf{Proof-status language.}  ``The last two steps are exact
finite computations'' is replaced by ``all by hand; scripts 13--15
corroborate''.  Proposition 13.4 and Theorem 13.5 now carry the label
[P$^c$] and begin ``Assume Theorem 12.1'', making the dependence on
Target A explicit.
\item \textbf{Reproducibility.}  Script 13 certifies the group law of
$Q_3/Q_7$ deterministically (associativity on all $729^3$ triples) and
certifies the character table exactly (each rounded row is a common
eigenvector of all class-sum matrices, computed in $\Z[\zeta_3]$ with
integer arithmetic, and has the right norm); script 15 certifies that
$\lambda_0$ is a character of $\hat H$ on all (generator, element) pairs,
and its certificate now includes the multiplication table of $Q_1/Q_3$;
a new \texttt{verify\_certificate.py} re-reads the certificate and
checks, without the cover, that the tables define an action of a group
of order $27$ with no fixed points and the listed orbits.  No sampled
checks remain in scripts 13--15.  The scripts used for this draft are
those of commit \texttt{67e6439} in the repository, and Appendix B says so.
\item \textbf{Prime degree.}  ``Every element of $D\setminus F$ generates
a maximal subfield'' is restricted to prime degree.
\end{enumerate}

\section*{Material on levels (Section 2.1, new paragraph)}
Level as the integer form of Moy--Prasad depth ($m/r$ for $D^\times$),
preservation of depth under Jacquet--Langlands, the finite groups
$D^\times/F^\times K^{m+1}$ and the count of representations of level
$\le2$ in the running example ($5+4+4$, dimensions $1,1,1,1,2;4;6$),
Howe--Corwin admissible pairs in the tame case $p\nmid r$ and the
relation of level to the conductor of $\chi$, and the remark that levels
are unbounded.  Bushnell--Henniart added to the bibliography.

\section*{Future directions (new Section 15)}
Why level two is a natural stopping point; the two statements (B1),
(B2) per level and which failures would be obstructions; the level-three
setup ($H_3$ of order $3^{11}$, Clifford theory over $Q_5/Q_{10}$, the
certificate format); two ideas toward an inductive mechanism (an
orbit-theoretic description in the spirit of Howe's Kirillov theory;
inductive freeness via the central subgroups $\hat A_m$, which are
central in $\hat H_m$ for $m\ge3$ with no hypothesis); what remains for
Target C; and the duality construction.

\section*{Numbering}
The new lemma shifts two numbers: the level-two theorem is now 13.8 and
freeness 13.9; the new sections shift the interface section to 16.
Labels unchanged; the scripts' README updated.
\end{document}
